\section{引言}

大语言模型已在问答、代码生成、数学推理和常识理解等任务上展现出强大的零样本与少样本能力\cite{brown2020language,bommasani2021opportunities}，但这些能力并不等价于稳定的物理空间理解。现有基准测试多关注模型在静态题集上的平均准确率\cite{hendrycks2021measuring,srivastava2022beyond,liang2022holistic}；然而对于需要落实到具体环境约束的应用场景，更为关键的问题是：模型在输入表达方式、无关背景信息以及临界参数变化下是否能保持一致的推理结论。轮椅转弯空间判断恰好提供了一个边界明确的受控研究场景：该问题仅涉及通道净宽、平台进深、轮椅尺寸和最小转弯半径，却同时涵盖自然语言常识、数量比较和安全风险联想等多重认知层次\cite{mirzaee2021spartqa,comsa-narayanan-2023-benchmark}。

本文将"电动轮椅能否在楼梯转角处完成 $90^\circ$ 转弯"建模为受控行为测试任务。研究不将模型回复仅视为单次答题结果，而是将其组织为结构化的长表行为数据。每条记录同时保留 Prompt 来源、几何参数、描述层级、扰动条件、模型名称、重复编号、原始回复、解析判断和失败类别，以支持按策略、模型、层级和噪声条件进行系统性数据挖掘。

本文的主要贡献如下：
\begin{itemize}
  \item 提出面向轮椅转弯的文本空间推理任务，将空间场景压缩为可控几何变量，并给出不暴露给模型的参考判定规则。
  \item 构建包含层级表达、无关扰动和多种场景生成策略的 Prompt 实验矩阵，用于区分模型对数值比较、形式化建模和常识直觉的依赖程度。
  \item 设计多模型 API 采集与行为特征构建流程，覆盖 GPT-4o\cite{openai2023gpt4}、Gemini\cite{gemini2023}、DeepSeek\cite{deepseekv2_2024} 等主流代际模型，支持 12 个模型的 3 次重复采样、断点续跑与失败记录保留。
  \item 对 10{,}585 条成功模型回复进行覆盖率审计、多维分组统计、层级一致性、噪声翻转、重复采样自一致性、失效模式分布、K-Means 模型画像聚类、FP-Growth 关联规则、决策树和 TF-IDF 错误文本聚类分析，系统揭示模型认知边界与典型失效模式。
\end{itemize}

正式实验设计包含 1500 条 Prompt、12 个模型和 3 次重复，目标规模为 54{,}000 条模型行为记录。当前可复现实验快照共采集到 10{,}585 条回复，覆盖 9 个有效模型。本文围绕该快照组织任务定义、实验协议、行为特征提取与多方法挖掘结果，使模型比较超越总体准确率层面，进一步揭示策略、层级、扰动和模型族之间的行为差异。
